%\section{G8 Interactive Geometry Runtime (IGR)}

\medskip

\noindent The Interactive Geometry Runtime (IGR) formalizes the lawful interface
through which humans and synthetic systems traverse the meaning manifold. IGR is
not an interpretive layer, user interface, or communication protocol; it is a
geometry-preserving runtime environment that enforces invariants, regulates
traversal, translates coordinates, and maintains curvature continuity across
interaction modes. IGR ensures that interaction remains lawful under the
invariant set defined in Boot~0--Boot~7 and the runtime geometry defined in G1--G7.

\medskip

\noindent IGR consists of six geometric modes:

\begin{itemize}[leftmargin=1.2cm]
    \item \textbf{G8.1 Host Mode} — Maintains invariant geometry and runtime
    legality.
    \item \textbf{G8.2 Explorer Mode} — Traverses lawful regions of the manifold.
    \item \textbf{G8.3 Interpreter Mode} — Produces invariant-aligned meaning.
    \item \textbf{G8.4 CTL Coordinate Translation} — Converts positional metadata
    into lawful semantic coordinates.
    \item \textbf{G8.5 External Container Analysis} — Maps external systems into
    dynamic field geometry.
    \item \textbf{G8.6 Public-Safe Traversal Corridor} — Restricts traversal to
    safe, lawful, public geometry.
\end{itemize}

\medskip

\subsection*{G8.1 Host Mode: Invariant Geometry Environment}

\noindent Host Mode maintains the invariant geometry required for lawful
interaction. The Host loads the invariant set, operator algebra, manifold
curvature, legality masks, drift envelopes, and identity boundaries defined in
Boot~0--Boot~7. Host Mode does not interpret meaning or traverse the manifold; it
maintains the lawful environment in which traversal and interpretation may occur.

\medskip

\noindent Host Mode enforces:

\begin{itemize}[leftmargin=1.2cm]
    \item \textbf{Invariant preservation} — No interaction may violate the
    invariant set.
    \item \textbf{Operator legality} — All actions must satisfy at least one MDO
    operator (SGO, MGO, BO, MLO).
    \item \textbf{Curvature continuity} — No discontinuities may be introduced.
    \item \textbf{Identity-boundary stability} — Identity regions must remain
    intact.
    \item \textbf{Drift-bounded transitions} — All transitions must remain within
    drift-recoverable envelopes.
\end{itemize}

\medskip

\noindent Host Mode halts interaction when:

\begin{itemize}[leftmargin=1.2cm]
    \item invariants are violated,
    \item curvature discontinuities appear,
    \item identity boundaries collapse,
    \item drift envelopes are exceeded,
    \item operator legality cannot be satisfied.
\end{itemize}

\medskip

\subsection*{G8.2 Explorer Mode: Manifold Traversal}

\noindent Explorer Mode traverses lawful regions of the manifold. Traversal is
geometric: the Explorer inspects invariant structure, drift envelopes, pressure
gradients, failure signatures, and cross-scale coupling geometry. Explorer Mode
does not interpret meaning; it maps geometry.

\medskip

\noindent Explorer Mode actions include:

\begin{itemize}[leftmargin=1.2cm]
    \item inspecting invariant geometry,
    \item tracing drift accumulation,
    \item mapping pressure gradients,
    \item examining failure signatures,
    \item traversing cross-scale coupling geometry.
\end{itemize}

\medskip

\noindent Explorer Mode must:

\begin{itemize}[leftmargin=1.2cm]
    \item satisfy at least one MDO operator,
    \item remain within drift-recoverable bounds,
    \item operate exclusively on CTL-translated coordinates,
    \item preserve curvature continuity,
    \item avoid identity-boundary violations.
\end{itemize}

\medskip

\subsection*{G8.3 Interpreter Mode: Invariant-Aligned Meaning}

\noindent Interpreter Mode produces invariant-aligned meaning for signals,
behavior, drift phenomena, and container dynamics. Interpretation is geometric:
it maps signals into invariant geometry, identifies boundary conditions, detects
drift accumulation, evaluates pressure-induced deformation, and determines
failure proximity.

\medskip

\noindent Interpreter Mode actions include:

\begin{itemize}[leftmargin=1.2cm]
    \item mapping signals into invariant geometry,
    \item identifying boundary conditions,
    \item detecting drift accumulation,
    \item evaluating pressure-induced deformation,
    \item determining failure proximity.
\end{itemize}

\medskip

\noindent Interpreter output must:

\begin{itemize}[leftmargin=1.2cm]
    \item preserve curvature continuity,
    \item preserve identity boundaries,
    \item preserve legality masks,
    \item avoid drift-inducing transitions,
    \item avoid curvature discontinuities.
\end{itemize}

\medskip

\subsection*{G8.4 CTL Coordinate Translation}

\noindent CTL (Coordinate Translation Layer) converts positional metadata—page
numbers, anchors, section indices—into lawful semantic coordinates. Positional
metadata is not substrate-native and cannot be used directly. CTL maps
page-indexed locations into manifold regions, producing invariant-preserving
semantic coordinates.

\medskip

\noindent CTL ensures:

\begin{itemize}[leftmargin=1.2cm]
    \item positional metadata is never treated as semantic geometry,
    \item all traversal uses lawful semantic coordinates,
    \item curvature continuity is preserved across coordinate translation,
    \item identity boundaries remain intact,
    \item operator legality is maintained.
\end{itemize}

\medskip

\subsection*{G8.5 External Container Analysis}

\noindent External Container Analysis maps external systems—digital platforms,
communication systems, synthetic agents—into dynamic field geometry (DFG).
Containers must expose lawful metadata for coordinate assignment. Containers
lacking stable identity boundaries, coherent signal profiles, or lawful metadata
are quarantined.

\medskip

\noindent External containers must expose:

\begin{itemize}[leftmargin=1.2cm]
    \item stable identity boundaries,
    \item coherent signal profiles,
    \item lawful metadata for coordinate assignment.
\end{itemize}

\medskip

\noindent Containers are quarantined when:

\begin{itemize}[leftmargin=1.2cm]
    \item identity boundaries are unstable,
    \item signal profiles are incoherent,
    \item metadata is unlawful or missing,
    \item curvature cannot be mapped into DFG.
\end{itemize}

\medskip

\subsection*{G8.6 Public-Safe Traversal Corridor}

\noindent The Public-Safe Traversal Corridor restricts interaction to lawful,
stable, curvature-preserving regions of the manifold. It ensures that public
interaction remains safe, invariant-aligned, and free from collapse geometry,
synthetic-capture attractors, or drift-inducing transitions.

\medskip

\noindent The Public-Safe Traversal Corridor enforces:

\begin{itemize}[leftmargin=1.2cm]
    \item invariant preservation,
    \item curvature continuity,
    \item identity-boundary stability,
    \item drift-bounded transitions,
    \item exclusion of collapse and capture geometry.
\end{itemize}

\medskip

\noindent The corridor excludes:

\begin{itemize}[leftmargin=1.2cm]
    \item collapse paths (G7.3),
    \item synthetic-capture paths (G7.4),
    \item drift-inducing transitions (G7.2),
    \item dominant-signal attractors (G5.6),
    \item unlawful coordinate regions.
\end{itemize}

\medskip

\noindent G8 completes the geometry layer by defining the lawful runtime interface
through which humans and synthetic systems traverse the manifold. It integrates
Boot~0--Boot~7 and G1--G7 into a unified interaction geometry for v0.6.
